% GENERATED FILE. Edit canonical lessons, then run npm run generate:books.
% canonical-source-hash: fnv1a64:338255503d70703d
% canonical-lessons: GU-C277-mat, GU-C277-mane-lage-che, GU-C277-karan, GU-C277-karan-ke, GU-R277-opinion-and-reason

\chapter{Opinions and reasons}
\label{ch:gu-opinions-and-reasons}
\hlchaptermodality{277}

\begin{chapteropening}
\textbf{By the end of this chapter:} \emph{I can say what I think (mane lāge che ke …, mārā mate), ask someone's view, and give a reason with kāraṇ ke or -ne līdhe.}

Complete GU-R277-opinion-and-reason: give an opinion and the reason for it.
\end{chapteropening}

\section[\texorpdfstring{\gu{મત} (mat)}{મત (mat)}]{\gu{મત} (mat) — an opinion; a view}

\label{lesson:GU-C277-mat}



\begin{quote}
Before the new one: say \textquotedblleft{}I think that …\textquotedblright{}, the A2 way. (\textbf{\gu{વિચારું} \gu{છું} \gu{કે}}.) Then \textquotedblleft{}your\textquotedblright{}, with respect. (\textbf{\gu{તમારું}}.)
\end{quote}

\subsection*{You'll want to know: \gu{મત}}
\textbf{\gu{મત}} — \emph{mat} — \textquotedblleft{}an opinion; a view\textquotedblright{}. It is masculine: \textbf{\gu{તમારો} \gu{મત}}, \textquotedblleft{}your opinion\textquotedblright{}.

\begin{quote}
\textbf{\gu{તમારો} \gu{મત} \gu{શું} \gu{છે}?}
\end{quote}

\begin{quote}
\emph{tamāro mat śũ che?} — \textquotedblleft{}What is your opinion?\textquotedblright{}
\end{quote}

The answer often opens on \textbf{\gu{મારા} \gu{મતે}} (\emph{mārā mate}), \textquotedblleft{}in my opinion\textquotedblright{}: \textbf{\gu{મત}} takes the \textbf{-\gu{ે}} of \textbf{\gu{સવારે}}, and \textbf{\gu{મારો}} turns to \textbf{\gu{મારા}} before it.

\begin{quote}
\textbf{\gu{મારા} \gu{મતે} \gu{ગુજરાતી} \gu{અઘરી} \gu{નથી}.}
\end{quote}

\begin{quote}
\emph{mārā mate gujarātī agharī nathī.} — \textquotedblleft{}In my opinion Gujarati is not hard.\textquotedblright{}
\end{quote}

\begin{cousinweb}
\textbf{\gu{મત}} (\emph{mat}) is a tatsama: Sanskrit \emph{mata}, \textquotedblleft{}thought, held to be\textquotedblright{}, and as a noun \textquotedblleft{}an opinion\textquotedblright{} (Monier-Williams), taken into Gujarati as it was. It belongs with the Sanskrit verb \emph{man-}, \textquotedblleft{}to think\textquotedblright{}.
\end{cousinweb}

\subsection*{Your turn}
Say these aloud:

\begin{itemize}
\raggedright
  \item \emph{mat}, then \emph{tamāro mat}
  \item \textquotedblleft{}What is your opinion?\textquotedblright{} — \emph{tamāro mat śũ che?}
  \item \textquotedblleft{}In my opinion he is a good doctor\textquotedblright{} — \emph{mārā mate te sāro dāktar che}
\end{itemize}

\begin{tcolorbox}[breakable,colback=teal!4,colframe=teal!35!black,title={Before you move on}]
What does \textbf{\gu{મત}} mean? (\textquotedblleft{}An opinion\textquotedblright{}.) \textquotedblleft{}In my opinion\textquotedblright{}? (\textbf{\gu{મારા} \gu{મતે}}.) Why \textbf{\gu{તમારો}}, not \textbf{\gu{તમારું}}? (\textbf{\gu{મત}} is masculine.)
\end{tcolorbox}
\section[\texorpdfstring{\gu{મને} \gu{લાગે} \gu{છે} \gu{કે} … (mane lāge che ke …)}{મને લાગે છે કે … (mane lāge che ke …)}]{\gu{મને} \gu{લાગે} \gu{છે} \gu{કે} … (mane lāge che ke …) — I think that …}

\label{lesson:GU-C277-mane-lage-che}



\begin{quote}
Say \textquotedblleft{}in my opinion\textquotedblright{}. (\emph{mārā mate}.) Then the verb \textbf{\gu{લાગવું}}: \textquotedblleft{}to cost; to seem\textquotedblright{}. Its second sense carries this lesson.
\end{quote}

\begin{grammarlens}[title={Grammar Lens: it seems to me}]
\begin{quote}
\textbf{\gu{મને} \gu{લાગે} \gu{છે} \gu{કે} \gu{વરસાદ} \gu{પડશે}.}
\end{quote}

\begin{quote}
\emph{mane lāge che ke varsād paḍśe.} — \textquotedblleft{}I think it will rain.\textquotedblright{}
\end{quote}

Word for word: \textquotedblleft{}to me it seems that it will rain\textquotedblright{}. The thinker takes \textbf{-\gu{ને}}, as with \textbf{\gu{મને} \gu{તાવ} \gu{છે}}: \textbf{\gu{તમને}}, \textbf{\gu{તેને}}, \textbf{\gu{અરુણને}}. \textbf{\gu{કે}} hands over the thought, which keeps its own verb.

In the past, \textbf{\gu{લાગ્યું}}: \textbf{\gu{મને} \gu{લાગ્યું} \gu{કે} \gu{તે} \gu{નહીં} \gu{આવે}} (\emph{mane lāgyũ ke te nahī̃ āve}), \textquotedblleft{}I thought he wouldn't come\textquotedblright{}. \textquotedblleft{}I don't think so\textquotedblright{} is \textbf{\gu{મને} \gu{નથી} \gu{લાગતું}} (\emph{mane nathī lāgtũ}).

\textbf{\gu{હું} \gu{વિચારું} \gu{છું} \gu{કે} …} is thinking something over. \textbf{\gu{મને} \gu{લાગે} \gu{છે} \gu{કે} …} gives an opinion.
\end{grammarlens}

\subsection*{You'll want to know: asking for a view}
\begin{quote}
\textbf{\gu{તમને} \gu{શું} \gu{લાગે} \gu{છે}?}
\end{quote}

\begin{quote}
\emph{tamne śũ lāge che?} — \textquotedblleft{}What do you think?\textquotedblright{}
\end{quote}

Two ways to ask: \textbf{\gu{તમને} \gu{શું} \gu{લાગે} \gu{છે}?} and \textbf{\gu{તમારો} \gu{મત} \gu{શું} \gu{છે}?} Two ways to answer: \textbf{\gu{મને} \gu{લાગે} \gu{છે} \gu{કે} …} and \textbf{\gu{મારા} \gu{મતે} …}: \textbf{\gu{મારા} \gu{મતે} \gu{ટ્રેન} \gu{મોડી} \gu{આવશે}} (\emph{mārā mate ṭren moḍī āvśe}).

\subsection*{Your turn}
Say these aloud:

\begin{itemize}
\raggedright
  \item \textquotedblleft{}I thought he wouldn't come\textquotedblright{} — \emph{mane lāgyũ ke te nahī̃ āve}
  \item \textquotedblleft{}I don't think so\textquotedblright{} — \emph{mane nathī lāgtũ}
\end{itemize}

\begin{tcolorbox}[breakable,colback=teal!4,colframe=teal!35!black,title={Before you move on}]
\textquotedblleft{}I think it will rain\textquotedblright{}? (\textbf{\gu{મને} \gu{લાગે} \gu{છે} \gu{કે} \gu{વરસાદ} \gu{પડશે}}.) \textquotedblleft{}What do you think?\textquotedblright{} (\textbf{\gu{તમને} \gu{શું} \gu{લાગે} \gu{છે}?}) Who takes \textbf{-\gu{ને}} in it? (The thinker.)
\end{tcolorbox}
\section[\texorpdfstring{\gu{કારણ} (kāraṇ)}{કારણ (kāraṇ)}]{\gu{કારણ} (kāraṇ) — a reason; a cause}

\label{lesson:GU-C277-karan}



\begin{quote}
Before the new one: ask \textquotedblleft{}why?\textquotedblright{} and answer \textquotedblleft{}because …\textquotedblright{}. (\textbf{\gu{કેમ}?} … \textbf{\gu{કેમકે}} ….) Then \textquotedblleft{}I think that …\textquotedblright{}. (\emph{mane lāge che ke …})
\end{quote}

\subsection*{You'll want to know: \gu{કારણ}}
\textbf{\gu{કારણ}} — \emph{kāraṇ} — \textquotedblleft{}a reason; a cause\textquotedblright{}. It is neuter: \textbf{\gu{તેનું} \gu{કારણ}}, \textquotedblleft{}the reason for it\textquotedblright{}.

\begin{quote}
\textbf{\gu{તેનું} \gu{કારણ} \gu{શું} \gu{છે}?}
\end{quote}

\begin{quote}
\emph{tenũ kāraṇ śũ che?} — \textquotedblleft{}What is the reason for it?\textquotedblright{}
\end{quote}

\begin{quote}
\textbf{\gu{મારી} \gu{પાસે} \gu{બે} \gu{કારણ} \gu{છે}.}
\end{quote}

\begin{quote}
\emph{mārī pāse be kāraṇ che.} — \textquotedblleft{}I have two reasons.\textquotedblright{}
\end{quote}

\begin{cousinweb}
\textbf{\gu{કારણ}} (\emph{kāraṇ}) is a tatsama: Sanskrit \emph{kāraṇa}, \textquotedblleft{}cause\textquotedblright{} (CDIAL 3057), \textquotedblleft{}what makes a thing happen\textquotedblright{}, from the root \emph{kṛ-}, \textquotedblleft{}to do\textquotedblright{}, that also stands behind \textbf{\gu{કરવું}}.
\end{cousinweb}

\subsection*{Your turn}
Say these aloud:

\begin{itemize}
\raggedright
  \item \emph{kāraṇ}, then \emph{tenũ kāraṇ}
  \item \textquotedblleft{}What is the reason for it?\textquotedblright{} — \emph{tenũ kāraṇ śũ che?}
  \item \textquotedblleft{}I have two reasons\textquotedblright{} — \emph{mārī pāse be kāraṇ che}
\end{itemize}

\begin{tcolorbox}[breakable,colback=teal!4,colframe=teal!35!black,title={Before you move on}]
What does \textbf{\gu{કારણ}} mean? (\textquotedblleft{}A reason; a cause\textquotedblright{}.) Its gender? (Neuter: \textbf{\gu{તેનું} \gu{કારણ}}.)
\end{tcolorbox}
\section[\texorpdfstring{\gu{કારણ} \gu{કે} · -\gu{ને} \gu{લીધે} (kāraṇ ke · -ne līdhe)}{કારણ કે · -ને લીધે (kāraṇ ke · -ne līdhe)}]{\gu{કારણ} \gu{કે}, -\gu{ને} \gu{લીધે} (kāraṇ ke, -ne līdhe) — because, because of}

\label{lesson:GU-C277-karan-ke}



\begin{quote}
Say \textquotedblleft{}a reason\textquotedblright{}. (\emph{kāraṇ}.) Then \textquotedblleft{}late\textquotedblright{}. (\textbf{\gu{મોડું}}.)
\end{quote}

\begin{grammarlens}[title={Grammar Lens: \gu{કારણ} \gu{કે}, and a clause}]
\begin{quote}
\textbf{\gu{હું} \gu{મોડો} \gu{આવ્યો}, \gu{કારણ} \gu{કે} \gu{ટ્રેન} \gu{મોડી} \gu{હતી}.}
\end{quote}

\begin{quote}
\emph{hũ moḍo āvyo, kāraṇ ke ṭren moḍī hatī.} — \textquotedblleft{}I came late, because the train was late.\textquotedblright{}
\end{quote}

\textbf{\gu{કારણ} \gu{કે}} (\emph{kāraṇ ke}) works like \textbf{\gu{કેમકે}}: what happened first, then the reason, a whole clause with its own verb. \textbf{\gu{મોડો}} follows the man who came, \textbf{\gu{મોડી}} the train.
\end{grammarlens}

\subsection*{You'll want to know: -\gu{ને} \gu{લીધે}, and a noun}
\begin{quote}
\textbf{\gu{વરસાદને} \gu{લીધે} \gu{અમે} \gu{ઘરે} \gu{રહ્યા}.}
\end{quote}

\begin{quote}
\emph{varsādne līdhe ame ghare rahyā.} — \textquotedblleft{}Because of the rain, we stayed at home.\textquotedblright{}
\end{quote}

When the reason is a single thing, a noun, it takes \textbf{-\gu{ને} \gu{લીધે}} (\emph{-ne līdhe}), \textquotedblleft{}because of\textquotedblright{}, or \textbf{-\gu{ને} \gu{કારણે}} (\emph{-ne kāraṇe}), the same thing with the noun you have:

\noindent
\begin{tabularx}{\linewidth}{@{}>{\raggedright\arraybackslash}X>{\raggedright\arraybackslash}X@{}}
\toprule
\textbf{the reason is} & \textbf{say} \\
\midrule
a clause: the train was late & \textbf{…, \gu{કારણ} \gu{કે} \gu{ટ્રેન} \gu{મોડી} \gu{હતી}} \\
a noun: the rain & \textbf{\gu{વરસાદને} \gu{લીધે} …} \\
a noun: the crowd & \textbf{\gu{ભીડને} \gu{કારણે} …} \\
\bottomrule
\end{tabularx}

\subsection*{Your turn}
Say these aloud:

\begin{itemize}
\raggedright
  \item \textquotedblleft{}Because of the crowd, we were late\textquotedblright{} — \emph{bhīḍne kāraṇe amne moḍũ thayũ}
  \item \textquotedblleft{}She won't come, because she has a fever\textquotedblright{} — \emph{te nahī̃ āve, kāraṇ ke tene tāv che}
\end{itemize}

\begin{tcolorbox}[breakable,colback=teal!4,colframe=teal!35!black,title={Before you move on}]
\textquotedblleft{}Because the train was late\textquotedblright{}? (\textbf{\gu{કારણ} \gu{કે} \gu{ટ્રેન} \gu{મોડી} \gu{હતી}}.) \textquotedblleft{}Because of the rain\textquotedblright{}? (\textbf{\gu{વરસાદને} \gu{લીધે}}.) Which one takes a noun? (\textbf{-\gu{ને} \gu{લીધે}}.)
\end{tcolorbox}
\section[(review)]{Review — opinions and reasons}

\label{lesson:GU-R277-opinion-and-reason}



\begin{quote}
Say \textquotedblleft{}I went home and ate\textquotedblright{}. (\emph{hũ ghare jaīne jamyo}.) Then \textquotedblleft{}first we drank tea, then we went to the market\textquotedblright{}. (\emph{pahelā̃ ame cā pīdhī, pachī ame bajār gayā}.)
\end{quote}

\subsection*{Your turn}
Say these aloud:

\begin{itemize}
\raggedright
  \item What do you think? What is your opinion?
  \item I think it will rain. In my opinion the train will come late.
  \item I came late, because the train was late.
  \item Because of the rain, we stayed at home. Because of the rain, I won't go.
  \item What is the reason for it? I have two reasons.
\end{itemize}

\begin{tcolorbox}[breakable,colback=teal!4,colframe=teal!35!black,title={Before you move on}]
\textquotedblleft{}In my opinion\textquotedblright{}? (\textbf{\gu{મારા} \gu{મતે}}.) \textquotedblleft{}I don't think so\textquotedblright{}? (\textbf{\gu{મને} \gu{નથી} \gu{લાગતું}}.) Which takes a clause, \textbf{\gu{કારણ} \gu{કે}} or \textbf{-\gu{ને} \gu{લીધે}}? (\textbf{\gu{કારણ} \gu{કે}}.)
\end{tcolorbox}
